\section{Turn syntax fields into an ordered bit sequence}
The reconstructed block and selected levels are still arrays of numbers. A
bitstream requires a second representation: an ordered sequence of fields with
specified widths. The value five in a three-bit field is \texttt{101}; the same
value in a five-bit field is \texttt{00101}. The field width determines leading
zeros, so writing the value as a machine word cannot replace bit emission.

H.263 generally transmits a field's most significant bit first. In this chapter's
output buffer, the first transmitted bit occupies bit seven of the first byte,
the next occupies bit six, and so on. A new field continues at the next free bit;
it does not automatically begin at a byte boundary. Byte order of a native
machine word and transmission order of a bit field are separate conventions.
The writer must produce the same output bytes in either simulator byte order.

Let $p$ be the number of already committed bits. Its next byte index and bit
position are
\[
  b=\left\lfloor\frac p8\right\rfloor,\qquad t=7-(p\bmod8).
\]
For a right-aligned value $v$ with declared width $n$, source bit $n-1$ is
emitted first. Then emit source bit $n-2$, continuing down to zero. If the field
crosses a byte boundary, recompute the next destination byte and bit position.
A twenty-two-bit picture start code is one example of a field wider than a byte;
the writer itself does not need to know what that field means.

\section{Keep committed bits separate from unused buffer storage}
The writer context contains a buffer pointer, byte capacity and committed bit
count. Capacity describes allocated storage. The count describes the valid stream
prefix. A partly used last byte also contains bits that do not yet belong to the
stream. Reading or writing the entire capacity to a file would expose unrelated
storage after the valid prefix.

For this exercise, preserve every bit outside the appended field, including low
bits after a partly filled last byte. That makes the boundary explicit and lets
tests detect accidental whole-byte replacement. Before a file is written, the
encoder must apply the syntax's actual alignment or termination rule. It cannot
assume that unused bits happen to be zero. A normal encoder can initialize its
arena to zero, but initialization alone does not change the number of committed
bits or authorize arbitrary padding between fields.

The supplied observation arena starts with bytes \texttt{A5}, whose bits are
\texttt{10100101}. Consider a context with three committed bits, followed by a
nine-bit append of decimal 419, or \texttt{110100011}:
\begin{center}
  \tagpdfsetup{table/header-rows={1}}
  \begin{tabular}{lrrl}
    \textbf{State}        & \textbf{Committed bits} & \textbf{First bytes} & \textbf{Meaning}                  \\
    Before append         & 3                       & \texttt{A5 A5}       & Only prefix \texttt{101} is valid \\
    Append nine bits      & 12                      & \texttt{BA 35}       & Last four bits still unused       \\
    Append four zero bits & 16                      & \texttt{BA 30}       & Two complete valid bytes
  \end{tabular}
\end{center}
After the nine-bit append, byte one contains the valid prefix \texttt{0011}
and retained low bits \texttt{0101}. Those low bits are not part of the field.
Appending four zeros explicitly replaces them and advances the count to sixteen.
This example is a packing trace; it is not a complete H.263 header.

\section{Expose a bounded append operation}
Implement \texttt{h263\_write\_bits}. Pass an eight-byte-aligned mapped 24-byte
context in \texttt{\$a0}, an unsigned right-aligned field value in \texttt{\$a1}
and a full-width declared bit count in \texttt{\$a2}. Context fields are unsigned
native-endian doublewords:
\begin{center}
  \tagpdfsetup{table/header-rows={1}}
  \begin{tabular}{rll}
    \textbf{Offset} & \textbf{Field}      & \textbf{Range or role}         \\
    0               & Buffer pointer      & Byte-addressed output arena    \\
    8               & Capacity in bytes   & 0 through 1\,048\,576          \\
    16              & Committed bit count & 0 through eight times capacity
  \end{tabular}
\end{center}
The complete context and buffer extents must be non-wrapping. Positive-capacity
buffer storage is mapped, needs no alignment beyond a byte, and is disjoint from
the context. A zero-capacity buffer is never accessed. Preserve the buffer pointer
and capacity, all saved general registers and the incoming stack pointer. Provide
an aligned stack and private space if your implementation needs a frame.

Admit widths zero through 32. Require $0\leq v<2^n$; the high bits must be zero
rather than silently discarded. Width zero therefore requires value zero.
Validate every full-width field and the existing context before buffer accesses.
Malformed context or field returns status one and count zero, retaining context
and output without accessing buffer bytes. Context reads needed for validation
are permitted.

After metadata admission, compare the requested width with the remaining bit
capacity. Insufficient space returns status two and count zero, again with no
buffer access or state change. A valid empty append succeeds with count zero
and performs no buffer access or context write. It still validates the context.

On nonempty success, append exactly the declared bits, preserve all bits outside
that window, and update only the committed count after output is complete.
Return status zero and the appended width in \texttt{\$v1}. Access only destination
bytes touched by the field. A wider load or store reaching neighbouring bytes
is outside this contract, even if those bytes are mapped. No file operation is
performed by this function.

\section{Make space failure a preflight decision}
Let byte capacity be $c$ and used bits be $p$. First establish $p\leq8c$.
Then the remaining space is $8c-p$, and admission requires
\[
  n\leq8c-p.
\]
The exercise's capacity bound makes the multiplication by eight safe. Subtracting
a validated used count also avoids relying on an unchecked $p+n$ calculation.
Only after the whole field fits should the first read-modify-write occur.

This gives field-level atomicity for ordinary admission failures: a rejected
field leaves the prior stream usable, including its count and partial last byte.
It does not yet give atomicity to a multi-field structure. If a coefficient event
requires several appends, its encoder must preflight their total space or stage
the complete event. A failed fifth append cannot be repaired merely by restoring
the old count, because earlier appends may already have replaced unused bits.

\section{Write one byte fragment at a time}
Begin with the bit-at-a-time path, then combine adjacent bits within a byte.
If the next destination offset is $r=p\bmod8$, that byte has $8-r$ available
positions. Take the smaller of the remaining field width and $8-r$. Extract
that many next source bits, shift them into their destination positions, clear
exactly the corresponding old bits, then OR in the replacement fragment.

For a fragment of $k$ bits ending above $s$ low positions, its mask is
$(2^k-1)\ll s$. Clearing that mask before OR is essential when the new fragment
contains zeros. An OR-only writer can set old zero bits, but cannot replace old
ones with zero. The initial \texttt{A5} data deliberately exposes that error.

The byte-chunk path is still ordered scalar work. A field carries at most 32
bits and may touch five destination bytes when unaligned. MSA is useful in the
transform and sample paths; it is not a requirement for this serial append.
Compare complete function work, including admission and cursor updates, before
claiming that fewer loop iterations improve the whole encoder.
